\documentclass[10pt]{amsart}

\usepackage[T1]{fontenc}
\usepackage[utf8]{inputenc}
\usepackage{lmodern}
\usepackage{microtype}
\usepackage{amsmath,amssymb,mathtools}
\usepackage{enumitem}
\usepackage[hidelinks]{hyperref}

\allowdisplaybreaks[1]
\setlength{\emergencystretch}{2em}
\setlist[enumerate]{leftmargin=2.3em,itemsep=0.2em,topsep=0.35em}
\numberwithin{equation}{section}

\newcommand{\N}{\mathbb N}
\newcommand{\Z}{\mathbb Z}
\newcommand{\PP}{\mathbb P}
\newcommand{\Gsem}{\langle G\rangle}
\newcommand{\ind}{\mathbf 1}
\newcommand{\e}{\mathrm e}
\newcommand{\dens}{\operatorname{dens}}

\newtheorem{theorem}{Theorem}[section]
\newtheorem{lemma}[theorem]{Lemma}
\newtheorem{proposition}[theorem]{Proposition}
\theoremstyle{remark}
\newtheorem{remark}[theorem]{Remark}

\title[One prime residue class covering a tail]
{One residue class modulo each prime can\\
cover every sufficiently large integer}
\author{}
\date{}

\subjclass[2020]{Primary 11B25; Secondary 11N35, 11N37}
\keywords{covering congruences, prime moduli, upper-bound sieve,
Selberg--Delange method, multiplicative semigroups}

\begin{document}

\begin{abstract}
We give an affirmative solution to Erd\H{o}s Problem 279. For every fixed
integer \(k\ge 3\), one can choose a residue \(0\le r_p<p\) for each prime
\(p\) so that every sufficiently large integer \(n\) can be written as
\(n=r_p+tp\) for some prime \(p\) and integer \(t\ge k\). The proof was
generated by OpenAI's GPT-5.6-sol model. Its main ingredients are a large
hub prime, a multiplicative semigroup, and a deterministic sieve that assigns
fresh prime residue classes without destroying earlier coverage.
\end{abstract}

\maketitle

\section{Introduction}

Let \(\N=\{1,2,\ldots\}\) and let \(\PP\) denote the set of primes. For a
fixed integer \(k\ge3\), a residue assignment is a sequence
\[
  \boldsymbol r=(r_p)_{p\in\PP},\qquad 0\le r_p<p.
\]
It covers \(n\) at level \(k\) if
\[
  n=r_p+tp
\]
for some \(p\in\PP\) and some \(t\in\Z\), \(t\ge k\). The problem considered
here, also indexed as Erd\H{o}s problem 279, asks whether for every \(k\ge3\)
some assignment covers a complete tail of \(\N\).

\begin{theorem}\label{thm:main}
For every integer \(k\ge3\), there are a residue assignment
\(\boldsymbol r=(r_p)_{p\in\PP}\) and an integer \(N\) such that every
\(n\ge N\) has a representation
\[
  n=r_p+tp
\]
with \(p\) prime and \(t\ge k\).
\end{theorem}

The use of canonical representatives is part of the assertion.

\begin{lemma}[Maturity equivalence]\label{lem:maturity}
If \(0\le r_p<p\), then
\[
 n=r_p+tp\ \text{for some }t\in\Z,\ t\ge k
 \quad\Longleftrightarrow\quad
 n\equiv r_p\pmod p\ \text{and}\ kp\le n.
\]
\end{lemma}

\begin{proof}
Only the reverse implication needs comment. Congruence makes
\(t=(n-r_p)/p\) an integer. If \(t\le k-1\), then
\[
 n=r_p+tp\le(p-1)+(k-1)p=kp-1,
\]
contrary to \(kp\le n\).
\end{proof}

If \(k_1\le k_2\), a cover at level \(k_2\) is a cover at level \(k_1\).
Thus a construction for \(k=3\) alone would not prove Theorem~\ref{thm:main};
our construction treats an arbitrary fixed \(k\).

\subsection{Outline}

We shift the target by writing \(m=n-1\) and select shifted classes
\[
  a_p\equiv r_p-1\pmod p.
\]
Fix a large prime \(h>k\), set
\[
 G=\{p\in\PP:p\equiv1\pmod h\},
 \qquad
 \delta=\frac1{h-1},
\]
and leave the shifted zero class on every prime outside \(\{h\}\cup G\).
The hub class \(a_h=1\) covers every \(G\)-smooth \(m\) not divisible by
\(h\). Apart from isolated prime targets, the only difficult integers are
\[
 m=h^e u,\qquad e\ge1,\qquad u\in\Gsem,
\]
where \(\Gsem\) is the multiplicative semigroup generated by \(G\).

The primes in \(G\) are partitioned into a finite hub set \(H\), a periodic
pool \(T\) used to cover isolated prime targets, and a large pool \(S\) used
for multiscale cleanup. For
\[
 \rho=\sqrt{h/k},\qquad X_{j+1}=\rho X_j,
\]
the controller annuli are
\[
 R_j=S\cap\left(\frac{X_{j+1}}h,\frac{X_j}k\right].
\]
They are disjoint and adjacent. A prime \(p\in R_j\) is both mature for
every target \(m\in(X_j,X_{j+1}]\) and strictly greater than \(m/h\).

The analytic core is a deterministic estimate, uniform in arbitrary
previously selected nonzero classes on the small \(G\)-primes. It bounds the
number of hard survivors at scale \(X\) by
\[
 \bigl(C_{\mathrm{sv}}+o(1)\bigr)
 \delta A_\delta v^\delta(\rho-1)\frac X{\log X},
 \qquad
 A_\delta=\frac{\e^{-\gamma\delta}}{\Gamma(\delta)},
\]
where \(C_{\mathrm{sv}}\) is absolute and \(v\) is fixed. Since
\[
 A_\delta<\delta,\qquad
 \frac{\rho-1}{1/k-\rho/h}=\sqrt{kh},
\]
this is smaller than the controller reservoir when \(h\) is large.

Permanence is supplied by a causal backup lemma. Changing a controller
\(p\in G\) from the provisional class \(0\) to a nonzero class can only
remove a zero-class witness \(cp\) below the current frontier with \(c<h\).
When \(c=1\), the hub \(h\) is a backup; when \(c>1\), a prime divisor of
\(c\) is a permanent outside zero-class backup.

\subsection{Analytic inputs}

We use Fouvry and Tenenbaum's Bombieri--Vinogradov theorem for
multiplicative functions \cite[Definition 1.1 and Theorem 1.8]{FouvryTenenbaum},
the fixed-function Selberg--Delange estimate of de la Bret\`eche and
Tenenbaum \cite[Theorem 2.1]{BretècheTenenbaum}, and the fundamental lemma
for upper beta-sieve weights \cite[Lemma 6.11]{FriedlanderIwaniec}, in the
convenient formulation \cite[(2.5)--(2.6), Theorem 2.2]{KurlbergEtAl}.
All hypotheses and constant dependencies are checked explicitly.

\section{A deterministic sieve for the hard semigroup}

Throughout this section, \(h\ge5\) is a fixed prime,
\[
 G=\{p\in\PP:p\equiv1\pmod h\},
 \qquad
 \delta=\frac1{h-1},
\]
and
\[
 f(u)=\ind_{u\in\Gsem}.
\]
The function \(f\) is completely multiplicative and takes values in
\(\{0,1\}\).

\begin{proposition}[Deterministic old-prime sieve]\label{prop:sieve}
Fix \(\rho>1\). There is an absolute constant \(C_{\mathrm{sv}}\) with the
following property. Choose a sufficiently large fixed integer \(v\). For
every \(X\), assign an arbitrary nonzero class \(a_p\bmod p\) to each
\[
 p\in G,\qquad p\le z:=X^{1/v}.
\]
Uniformly in all these choices, the number of integers
\[
 m=h^e u\in(X,\rho X],\qquad e\ge1,\qquad u\in\Gsem,
\]
which avoid every assigned class is at most
\begin{equation}\label{eq:main-sieve}
 \bigl(C_{\mathrm{sv}}+o_{h,\rho,v}(1)\bigr)
 \delta A_\delta v^\delta(\rho-1)\frac X{\log X},
 \qquad
 A_\delta=\frac{\e^{-\gamma\delta}}{\Gamma(\delta)}.
\end{equation}
\end{proposition}

\subsection{Translated progression estimates}

Fix \(e\ge1\). For \(p\in G\), define the nonzero class
\begin{equation}\label{eq:bpe}
 b_{p,e}\equiv a_p h^{-e}\pmod p.
\end{equation}
For a squarefree product \(d\) of primes in \(G\), CRT gives a reduced class
\(b_{d,e}\bmod d\) satisfying the local congruences. Since \((d,h)=1\),
choose an integer representative with
\begin{equation}\label{eq:btilde}
 \widetilde b_{d,e}\equiv b_{d,e}\pmod d,
 \qquad
 \widetilde b_{d,e}\equiv1\pmod h.
\end{equation}
Then \((\widetilde b_{d,e},dh)=1\), while the progression modulo \(d\) is
unchanged.

For \(I=(Y,\rho Y]\), write
\[
 F(I)=\sum_{u\in I}f(u),
 \qquad
 F_d(I)=\sum_{\substack{u\in I\\(u,d)=1}}f(u),
 \qquad
 A_d(I,b)=\sum_{\substack{u\in I\\u\equiv b\;(\mathrm{mod}\ d)}}f(u).
\]
The class \(\mathcal F(D,K)\) in \cite[Definition 1.1]{FouvryTenenbaum}
consists of multiplicative functions \(g\) for which:
\begin{enumerate}[label=\textup{(\roman*)}]
\item there are \(2=\Upsilon_1<\Upsilon_2<\cdots\) satisfying
\[
 \frac{\Upsilon_{j+1}}{\Upsilon_j}
 \ge1+\frac1{(\log(2\Upsilon_j))^K};
\]
\item \(g(p)=g(p')\) whenever \(p,p'\) lie in the same corresponding
interval and \(p\equiv p'\pmod D\);
\item \(\lvert g(n)\rvert\le\tau_K(n)\).
\end{enumerate}
Our \(f\) lies in \(\mathcal F(h,1)\): its value on primes depends only on
the class modulo \(h\), and \(\lvert f\rvert\le1=\tau_1\).

For every \(A>0\), \cite[Theorem 1.8]{FouvryTenenbaum} supplies
\(B=B(A,1)\) such that
\begin{equation}\label{eq:BV}
 \sum_{q\le\sqrt{x}/(\log(3x))^B}
 \max_{(b,qh)=1}\lvert\Delta_f(x;q,h,b)\rvert
 \ll_{h,A}\frac{x}{(\log(3x))^A}.
\end{equation}
When \((q,h)=1\), the discrepancy is
\[
 \Delta_f(x;q,b)
 =
 \sum_{\substack{u\le x\\u\equiv b\;(\mathrm{mod}\ q)}}f(u)
 -\frac1{\varphi(q)}
 \sum_{\substack{u\le x\\(u,q)=1}}f(u).
\]
Define
\begin{equation}\label{eq:DY}
 D_Y=\frac{\sqrt Y}{(\log(3\rho Y))^B}.
\end{equation}
This is within the range in \eqref{eq:BV} at both \(x=Y\) and \(x=\rho Y\).
Applying \eqref{eq:BV} at the endpoints, using \eqref{eq:btilde}, and
restricting to squarefree \(G\)-supported moduli gives
\begin{equation}\label{eq:translated-BV}
 \sum_{d\le D_Y}\max_{(b,dh)=1}
 \left|A_d(I,b)-\frac{F_d(I)}{\varphi(d)}\right|
 \ll_{h,A,\rho}\frac Y{(\log Y)^A}.
\end{equation}
The maximum makes this uniform in every correlated vector of old residue
classes.

\subsection{A uniform coprime mean by exact M\"obius inversion}

The prime number theorem in the fixed progression \(1\bmod h\) gives
\[
 \sum_{p\le x}f(p)\log p=\delta x+O_h(x/\log^2x).
\]
For any fixed \(1/2<\sigma_0<1\),
\[
 \sum_p\left\{
 \frac{f(p)^2}{p^{2\sigma_0}}
 +\sum_{\nu\ge2}\frac{f(p^\nu)}{p^{\nu\sigma_0}}
 \right\}<\infty.
\]
Thus \cite[Theorem 2.1]{BretècheTenenbaum} applies with \(A=2\) and
\(r=\delta\). In that theorem \(\delta_{1,A}\) is Kronecker's delta, so
\(\delta_{1,2}=0\). Hence
\begin{equation}\label{eq:M-asymp}
 M(x):=\sum_{u\le x}f(u)
 =\frac{K_h}{\Gamma(\delta)}x(\log x)^{\delta-1}
 \{1+O_h(1/\log x)\},
\end{equation}
where
\[
 K_h=\lim_{z\to\infty}K_h(z)>0,
\]
\begin{equation}\label{eq:Kh}
 K_h(z)=
 \prod_{p\le z}(1-1/p)^\delta
 \prod_{\substack{p\le z\\p\in G}}(1-1/p)^{-1}.
\end{equation}
No parameter-uniform contour theorem is required. If \(d\) is squarefree
and supported on \(G\), exact M\"obius inversion gives
\begin{align}
 M_d(x)
 &:=
 \sum_{\substack{u\le x\\(u,d)=1}}f(u) \notag\\
 &=
 \sum_{r\mid d}\mu(r)\sum_{\substack{u\le x\\r\mid u}}f(u)
 =
 \sum_{r\mid d}\mu(r)M(x/r).
\label{eq:mobius}
\end{align}
The final identity holds because \(r\in\Gsem\) and \(f(rv)=f(v)\).

\begin{lemma}[Local density]\label{lem:local-density}
Uniformly for squarefree \(d\le\sqrt Y\) supported on \(G\),
\begin{equation}\label{eq:local-density}
 \frac{F_d((Y,\rho Y])}{\varphi(d)}
 =
 \frac{F((Y,\rho Y])}{d}
 \left\{1+O_{h,\rho}\left(
 \frac{(\log\log(3Y))^3}{\log Y}\right)\right\}.
\end{equation}
\end{lemma}

\begin{proof}
For \(x\in\{Y,\rho Y\}\) and \(r\mid d\), one has \(x/r\ge\sqrt Y\).
Inserting \eqref{eq:M-asymp} into \eqref{eq:mobius}, using
\[
 \left(1-\frac{\log r}{\log x}\right)^{\delta-1}
 =1+O_h\left(\frac{\log r}{\log x}\right)
 \qquad(r\le\sqrt x),
\]
and then applying
\[
 \sum_{r\mid d}\frac1r=\prod_{p\mid d}(1+1/p),
 \qquad
 \sum_{r\mid d}\frac{\log r}{r}
 =
 \prod_{p\mid d}(1+1/p)
 \sum_{p\mid d}\frac{\log p}{p+1},
\]
gives
\begin{equation}\label{eq:Md}
 M_d(x)=\frac{\varphi(d)}dM(x)
 \left\{1+O_h\left(\frac{R(d)}{\log x}\right)\right\},
\end{equation}
where
\[
 R(d)=
 \prod_{p\mid d}\frac{p+1}{p-1}
 \left(1+\sum_{p\mid d}\frac{\log p}{p+1}\right).
\]
The maximal-order estimate \(d/\varphi(d)\ll\log\log(3d)\) gives
\[
 \prod_{p\mid d}\frac{p+1}{p-1}
 \le\left(\frac d{\varphi(d)}\right)^2
 \ll(\log\log(3Y))^2.
\]
Also
\[
 1+\sum_{p\mid d}\frac{\log p}{p+1}
 \ll\log\log(3Y).
\]
Indeed, split at \(Z=\log(3d)\); the contribution of \(p\le Z\) is
\(O(\log Z)\), and the contribution of \(p>Z\) is at most
\(Z^{-1}\sum_{p\mid d}\log p=O(1)\). Thus
\(R(d)\ll(\log\log(3Y))^3\).

Subtract \eqref{eq:Md} at \(Y\) and at \(\rho Y\). Because \(\rho>1\) is
fixed,
\[
 M(\rho Y)-M(Y)\asymp_{h,\rho}Y(\log Y)^{\delta-1},
\]
so both endpoint errors have the asserted relative size.
\end{proof}

\subsection{Translated upper sieve}

Let
\[
 P(z)=\prod_{\substack{p\le z\\p\in G}}p.
\]
Use upper beta-sieve weights \(\lambda_d^+\) of level \(D_Y\). Their
required properties are
\begin{equation}\label{eq:sieve-support}
 |\lambda_d^+|\le1,\qquad
 \lambda_d^+=0\ \text{unless }d\mid P(z),\ d\text{ squarefree},\ d\le D_Y,
\end{equation}
and
\begin{equation}\label{eq:sieve-majorant}
 \ind_{L=1}\le\sum_{d\mid L}\lambda_d^+
 \qquad(L\mid P(z)).
\end{equation}
For \(g(d)=1/d\), the sieve has dimension at most one, since
\begin{equation}\label{eq:sieve-dimension}
 \prod_{\substack{w\le p<z\\p\in G}}(1-1/p)^{-1}
 \le
 \prod_{w\le p<z}(1-1/p)^{-1}
 \le
 \frac{\log z}{\log w}\{1+O(1/\log w)\}.
\end{equation}
The constant is absolute. The fundamental lemma
\cite[Lemma 6.11]{FriedlanderIwaniec} therefore supplies absolute constants
\(s_0\) and \(C_{\mathrm{sv}}\) such that, if \(D_Y\ge z^{s_0}\),
\begin{equation}\label{eq:fundamental-lemma}
 \sum_{d\mid P(z)}\frac{\lambda_d^+}{d}
 \le C_{\mathrm{sv}}
 \prod_{\substack{p\le z\\p\in G}}(1-1/p).
\end{equation}
More precisely, one may fix a beta-sieve parameter \(\beta\ge5\); for
\(s=\log D_Y/\log z\ge\beta+1\), the left side equals the product times
\(1+O(s^{-s/2})\) \cite[Theorem 2.2]{KurlbergEtAl}.

For \(u\in I\), define
\[
 L_e(u)=
 \prod_{\substack{p\le z,\ p\in G\\u\equiv b_{p,e}\;(\mathrm{mod}\ p)}}p.
\]
Survival is exactly \(L_e(u)=1\), and \(d\mid L_e(u)\) is exactly the CRT
fibre \(u\equiv b_{d,e}\pmod d\). Thus \eqref{eq:sieve-majorant} gives
\begin{align}
 S_e(I)
 &:=
 \sum_{\substack{u\in I\\
 u\not\equiv b_{p,e}\;(\mathrm{mod}\ p)\ (p\le z,\ p\in G)}}f(u)
 \notag\\
 &\le
 \sum_{d\mid P(z)}\lambda_d^+A_d(I,b_{d,e}).
\label{eq:Se-majorant}
\end{align}
Set
\[
 V_h(z)=\prod_{\substack{p\le z\\p\in G}}(1-1/p),
 \qquad
 \eta_Y=\frac{(\log\log(3Y))^3}{\log Y}.
\]
Combining \eqref{eq:translated-BV}, \eqref{eq:local-density}, and
\eqref{eq:Se-majorant} gives
\begin{align}
 S_e(I)
 &\le
 F(I)\sum_d\frac{\lambda_d^+}{d}
 +O_{h,\rho}\left(F(I)\eta_Y\sum_d\frac{|\lambda_d^+|}{d}\right)
 +O_{h,A,\rho}\left(\frac Y{(\log Y)^A}\right).
\label{eq:Se-estimate}
\end{align}
Now
\[
 \sum_d\frac{|\lambda_d^+|}{d}
 \le
 \prod_{\substack{p\le z\\p\in G}}(1+1/p)
 \ll_h(\log z)^\delta,
\]
whereas \(V_h(z)\asymp_h(\log z)^{-\delta}\). Relative to
\(F(I)V_h(z)\), the middle error in \eqref{eq:Se-estimate} is
\[
 \ll_h
 \frac{(\log\log(3Y))^3}{\log Y}(\log z)^{2\delta}
 =o(1),
\]
because \(h\ge5\) gives \(2\delta<1\). Taking \(A>2\) makes the final error
negligible. Hence
\begin{equation}\label{eq:Se-final}
 S_e(I)\le\{C_{\mathrm{sv}}+o_{h,\rho}(1)\}F(I)V_h(z).
\end{equation}

\subsection{Exact constant and summation over the layers}

The constant \(K_h\) in \eqref{eq:M-asymp} also determines the
fixed-progression Mertens formula. Indeed, \eqref{eq:Kh} gives
\[
 V_h(z)^{-1}
 =
 K_h(z)\prod_{p\le z}(1-1/p)^{-\delta}
 \sim K_h\e^{\gamma\delta}(\log z)^\delta
\]
by the ordinary Mertens product. Equivalently,
\begin{equation}\label{eq:Vh}
 V_h(z)\sim\e^{-\gamma\delta}K_h^{-1}(\log z)^{-\delta}.
\end{equation}
Therefore
\begin{align}
 F((Y,\rho Y])V_h(z)
 &\sim
 \frac{\e^{-\gamma\delta}}{\Gamma(\delta)}
 (\rho-1)\frac Y{\log Y}
 \left(\frac{\log Y}{\log z}\right)^\delta
 \notag\\
 &=
 A_\delta(\rho-1)\frac Y{\log Y}
 \left(\frac{\log Y}{\log z}\right)^\delta.
\label{eq:constant-cancel}
\end{align}
Thus the \(h\)-dependent constant \(K_h\) cancels exactly.

Take \(z=X^{1/v}\). For \(m=h^eu\in(X,\rho X]\), put \(Y_e=X/h^e\).
If \(Y_e\ge\sqrt X\), then
\begin{equation}\label{eq:sieve-level-ratio}
 \frac{\log D_{Y_e}}{\log z}
 =
 \frac{\tfrac12\log Y_e-B\log\log(3\rho Y_e)}
 {(1/v)\log X}
 \ge\frac v4-o(1).
\end{equation}
Choose \(v\) so that \(v/4>s_0\). Since
\((\log Y_e/\log z)^\delta\le v^\delta\), equations
\eqref{eq:Se-final} and \eqref{eq:constant-cancel} give
\begin{equation}\label{eq:Se-layer}
 S_e\le\{C_{\mathrm{sv}}+o(1)\}
 A_\delta v^\delta(\rho-1)
 \frac{X/h^e}{\log(X/h^e)}.
\end{equation}
All error terms used above are uniform for \(X^{1/2}\le Y_e\le X\).
To justify the next asymptotic, first fix \(E\) and let \(X\to\infty\) in
the terms \(e\le E\). For the remaining terms,
\(\log(X/h^e)\ge\frac12\log X\), so after division by \(X/\log X\)
their contribution is at most \(2\sum_{e>E}h^{-e}\). Now let
\(E\to\infty\). Thus
\begin{equation}\label{eq:layer-sum}
 \sum_{\substack{e\ge1\\X/h^e\ge\sqrt X}}
 \frac{X/h^e}{\log(X/h^e)}
 \sim
 \frac X{\log X}\sum_{e\ge1}h^{-e}
 =
 \delta\frac X{\log X}.
\end{equation}
For the complementary layers the trivial estimate is
\begin{equation}\label{eq:small-layers}
 \sum_{\substack{e:Y_e<\sqrt X\\\rho Y_e\ge1}}
 \bigl(|(Y_e,\rho Y_e]\cap\N|+1\bigr)
 \ll_{h,\rho}\sqrt X+\log X
 =o(X/\log X).
\end{equation}
All remaining intervals contain no positive integer. Equations
\eqref{eq:Se-layer}--\eqref{eq:small-layers} prove
Proposition~\ref{prop:sieve}.

\begin{lemma}\label{lem:Adelta}
For \(0<\delta<1\), one has \(A_\delta<\delta\).
\end{lemma}

\begin{proof}
The functional equation and Weierstrass product for the gamma function give
\[
 \frac{A_\delta}{\delta}
 =
 \frac{\e^{-\gamma\delta}}{\Gamma(1+\delta)}
 =
 \prod_{n\ge1}\left(1+\frac\delta n\right)\e^{-\delta/n}
 <1,
\]
because \(\log(1+x)<x\) for \(x>0\).
\end{proof}

\section{Prime targets and the controller reservoir}

Fix an arbitrary integer \(k\ge3\). Choose the absolute integer \(v\)
required in Proposition~\ref{prop:sieve}. We next choose a prime \(h>k\),
sufficiently large that
\begin{equation}\label{eq:h-choice}
 2C_{\mathrm{sv}}v^\delta\sqrt{kh}\,A_\delta<1,
 \qquad
 \delta=(h-1)^{-1}.
\end{equation}
This is possible by Lemma~\ref{lem:Adelta}, since the left side is at most
\[
 \frac{2C_{\mathrm{sv}}v^\delta\sqrt{kh}}{h-1}\longrightarrow0
\]
as \(h\to\infty\) through the primes.

Set
\[
 \rho=\sqrt{h/k},
 \qquad
 G=\{p:p\equiv1\pmod h\}.
\]
In the shifted variable \(m=n-1\), begin with
\begin{equation}\label{eq:initial-classes}
 a_h=1,\qquad a_p=0\quad(p\notin\{h\}\cup G).
\end{equation}
Every class subsequently assigned on \(G\) will be nonzero.

\subsection{Covering isolated prime targets}

Choose a finite set \(H\subset G\), and set \(a_\ell=1\) for \(\ell\in H\).
A sufficiently large prime target \(q\) is already covered if
\[
 q\equiv1\pmod h
 \quad\text{or}\quad
 q\equiv1\pmod\ell\ \text{for some }\ell\in H.
\]
By the prime number theorem in the finitely many reduced CRT classes, the
complementary primes have density
\begin{equation}\label{eq:DeltaH}
 \Delta_H=
 \left(1-\frac1{h-1}\right)
 \prod_{\ell\in H}\left(1-\frac1{\ell-1}\right).
\end{equation}
The series \(\sum_{\ell\in G}1/\ell\) diverges, so \(\Delta_H\) can be made
arbitrarily small.

Enlarge \(H\) until a number \(\tau\) can satisfy
\begin{equation}\label{eq:tau-choice}
 k\Delta_H<\tau<h\Delta_H,
 \qquad
 \tau<\delta/2.
\end{equation}
We realize \(\tau\) as a periodic prime density. Choose a sufficiently large
auxiliary prime \(L\), coprime to \(h\prod_{\ell\in H}\ell\). Each pair of
reduced congruences
\[
 p\equiv1\pmod h,
 \qquad
 p\equiv c\pmod L
\]
defines a prime set of density \(1/((h-1)(L-1))\). Unions of such classes
have density mesh \(\delta/(L-1)\). Thus, for large \(L\), a union
\(T_0\subset G\) has density \(\tau\) satisfying \eqref{eq:tau-choice}. Put
\begin{equation}\label{eq:pools}
 T=T_0\setminus H,\qquad
 S=G\setminus(H\cup T),\qquad
 \sigma=\dens_{\PP}(S)=\delta-\tau>\delta/2.
\end{equation}

Let \(q_1<q_2<\cdots\) enumerate the complementary primes in
\eqref{eq:DeltaH}, and let \(p_1<p_2<\cdots\) enumerate \(T\). Inverting
the fixed-progression prime-counting asymptotics gives
\begin{equation}\label{eq:piqi}
 \frac{p_i}{q_i}\longrightarrow\frac{\Delta_H}{\tau}
 \in(1/h,1/k).
\end{equation}
For every sufficiently large \(i\), define
\begin{equation}\label{eq:prime-target-class}
 a_{p_i}\equiv q_i\pmod{p_i}.
\end{equation}
This class is nonzero and \(kp_i<q_i<q_i+1\). Assign \(a_p=1\) to the
finitely many unused early primes of \(T\), and absorb the finitely many
early \(q_i\) into the final threshold. Thus every sufficiently large shifted
prime target is covered.

\subsection{Disjoint controller annuli}

Choose a starting scale \(X_0\), to be enlarged below, and define
\begin{equation}\label{eq:annuli}
 X_j=X_0\rho^j,\qquad
 R_j=S\cap\left(\frac{X_{j+1}}h,\frac{X_j}k\right]
 \quad(j\ge0).
\end{equation}
Since \(\rho^2=h/k\), one has \(X_{j+2}/h=X_j/k\). The half-open annuli
\(R_j\) are pairwise disjoint and tile a tail of \(S\). The fixed-modulus
prime number theorem gives
\begin{equation}\label{eq:Rj-size}
 |R_j|\sim
 \sigma\left(\frac1k-\frac\rho h\right)\frac{X_j}{\log X_j}.
\end{equation}

Assign \(a_p=1\) to each \(S\)-prime \(p\le X_1/h\). At the beginning of
stage \(j\), every \(G\)-prime \(p\le X_j^{1/v}\) has a permanent nonzero
class. This is clear for \(H\cup T\). For \(S\), the initial segment and
earlier annuli suffice once \(X_0\) is large, because
\[
 hX_j^{1/v}<X_j.
\]
Call
\begin{equation}\label{eq:hard-target}
 m=h^eu,\qquad e\ge1,\qquad u\in\Gsem
\end{equation}
a hard target. Apply Proposition~\ref{prop:sieve} to
\((X_j,X_{j+1}]\), ignoring every more recent prime. The number of hard
targets not covered by old primes is at most
\begin{equation}\label{eq:survivors}
 \{C_{\mathrm{sv}}+o(1)\}
 \delta A_\delta v^\delta(\rho-1)\frac{X_j}{\log X_j}.
\end{equation}
The identity
\begin{equation}\label{eq:ratio-identity}
 \frac{\rho-1}{1/k-\rho/h}=\sqrt{kh}
\end{equation}
and \(\sigma>\delta/2\) show that the ratio of the leading coefficient in
\eqref{eq:survivors} to that in \eqref{eq:Rj-size} is at most
\[
 2C_{\mathrm{sv}}v^\delta\sqrt{kh}\,A_\delta<1.
\]
After fixing \(h\), choose \(X_0\) so large that this strict inequality beats
all error terms at every stage, all preceding asymptotics are valid, and
\begin{equation}\label{eq:X0-choice}
 X_1/h>kh.
\end{equation}

At stage \(j\), list the surviving hard targets and the primes of \(R_j\)
increasingly. Pair each survivor with a distinct prime \(p\in R_j\). If
\(p\) is paired with \(m\), put
\begin{equation}\label{eq:controller-class}
 a_p\equiv m\pmod p;
\end{equation}
give every unused prime in \(R_j\) the class \(a_p=1\).

\begin{lemma}[Controller validity]\label{lem:controller}
The class in \eqref{eq:controller-class} is nonzero, and \(p\) is mature at
the original target \(m+1\).
\end{lemma}

\begin{proof}
Since \(p>X_{j+1}/h\ge m/h\), divisibility \(p\mid m=h^eu\) would imply
\(m/p\ge h\), a contradiction. Hence \(m\not\equiv0\pmod p\). Moreover,
\[
 p\le X_j/k<m/k<(m+1)/k.
\]
\end{proof}

\subsection{Causal permanence}

We temporarily regard an unprocessed \(S\)-prime as having class zero.

\begin{lemma}[Causal backup]\label{lem:backup}
Let \(p\in R_j\). Changing \(a_p\) from \(0\) to an arbitrary nonzero
class cannot make any positive shifted integer \(x\le X_{j+1}\) uncovered,
provided all other assignments already used below that frontier are kept.
\end{lemma}

\begin{proof}
A point that could lose \(p\) as a zero-class witness has the form \(x=cp\).
Since \(p>X_{j+1}/h\), we have \(c<h\).

If \(c=1\), then \(p\equiv1\pmod h\), so the permanent hub class \(a_h=1\)
covers \(x=p\). Its original target \(p+1\) is mature because
\(kh<p<p+1\), by \eqref{eq:X0-choice}.

If \(c>1\), choose a prime divisor \(d\mid c\). Then \(d<h\), hence
\(d\notin\{h\}\cup G\). The permanent class \(a_d=0\) covers \(x=cp\),
and it is mature because
\[
 kd<kh<p<cp+1.
\]
\end{proof}

It follows inductively that every hard target covered at a stage remains
covered forever. The initial segment and adjacent annuli assign one nonzero
class to every prime in \(S\); there is no reselection.

\section{Completion of the cover}

\begin{proposition}\label{prop:completion}
With the assignment constructed in Section~3, every sufficiently large
shifted integer \(m\) is covered by a prime \(p\) satisfying
\[
 kp\le m+1.
\]
\end{proposition}

\begin{proof}
There are four cases.

\smallskip
\noindent\emph{Case 1: \(m\) is \(\{h\}\cup G\)-smooth and \(h\nmid m\).}
Every prime factor of \(m\) then lies in \(G\), hence is \(1\bmod h\).
Thus \(m\equiv1=a_h\pmod h\). For \(m+1\ge kh\), the hub is mature.

\smallskip
\noindent\emph{Case 2: \(m\) is \(\{h\}\cup G\)-smooth and \(h\mid m\).}
Then \(m\) has the unique hard form
\[
 m=h^eu,\qquad e\ge1,\qquad u\in\Gsem,
\]
and was permanently covered in its block.

\smallskip
\noindent\emph{Case 3: \(m\) has a prime factor
\(q\notin\{h\}\cup G\) with \(q\le m/k\).}
The permanent outside class \(a_q=0\) covers \(m\), and
\(kq\le m<m+1\).

\smallskip
\noindent\emph{Case 4: \(m\) is not \(\{h\}\cup G\)-smooth and every
outside prime factor exceeds \(m/k\).}
Assume \(m>k^2\). Since \(m\) is not \(\{h\}\cup G\)-smooth, it has at
least one outside prime factor. Two outside prime factors, counted with
multiplicity, would have product \(>m^2/k^2>m\), which is impossible.
Therefore there is exactly one outside prime factor \(q\), occurring to the
first power. The cofactor \(c=m/q\) is \(\{h\}\cup G\)-smooth and satisfies
\(c<k\). Every nontrivial generator of this semigroup is at least \(h>k\),
so \(c=1\). Hence \(m=q\) is prime, and it is covered by the prime-target
construction.

These cases exhaust all \(m>k^2\). The finitely many early prime-target
exceptions and the fixed maturity thresholds are absorbed by the phrase
``sufficiently large'' in the proposition.
\end{proof}

\begin{proof}[Proof of Theorem~\ref{thm:main}]
Let \(M_0\) exceed \(X_0\), \(k^2\), all finitely many early prime-target
exceptions, and every fixed maturity threshold. By
Proposition~\ref{prop:completion}, for every \(m\ge M_0\) there is a prime
\(p\) such that
\[
 m\equiv a_p\pmod p,\qquad kp\le m+1.
\]
For each prime \(p\), define \(r_p\) to be the unique element of
\(\{0,\ldots,p-1\}\) satisfying
\[
 r_p\equiv a_p+1\pmod p.
\]
Take \(N=M_0+1\). If \(n\ge N\), put \(m=n-1\); then
\[
 n\equiv r_p\pmod p,\qquad kp\le n.
\]
Lemma~\ref{lem:maturity} gives \(n=r_p+tp\) for an integer \(t\ge k\).
The assignment is fixed globally and uses exactly one canonical class
modulo every prime.
\end{proof}

\section{A compactness reformulation}

Although compactness is not needed in the construction, it clarifies the
global quantifier. Fix \(k\). For \(N\le M\), let \(F(k,N,M)\) be the finite
constraint problem that chooses one residue modulo every prime \(p\le M/k\)
and covers each \(n\in[N,M]\) using a prime \(p\le n/k\).

\begin{proposition}\label{prop:compactness}
The following are equivalent:
\begin{enumerate}[label=\textup{(\roman*)}]
\item one residue assignment covers every sufficiently large integer at
level \(k\);
\item there is a fixed \(N\) such that \(F(k,N,M)\) is satisfiable for every
\(M\ge N\).
\end{enumerate}
\end{proposition}

\begin{proof}
Let
\[
 \Omega=\prod_{p\in\PP}\{0,\ldots,p-1\},
\]
with finite discrete factors. It is compact. For each \(n\), the set
\[
 K_n=
 \{\boldsymbol r\in\Omega:
 \exists p\le n/k,\ r_p\equiv n\pmod p\}
\]
is clopen and depends on finitely many coordinates. A solution of
\(F(k,N,M)\) is equivalent, after arbitrary extension on unused coordinates,
to a point of
\[
 L_M=\bigcap_{n=N}^M K_n.
\]
If (ii) holds, the sets \(L_M\) are nonempty, closed, and nested.
Compactness gives
\[
 \bigcap_{M\ge N}L_M=\bigcap_{n\ge N}K_n\ne\varnothing,
\]
which is (i). The converse is immediate by restriction.
\end{proof}

\begin{remark}
Negating Proposition~\ref{prop:compactness} gives
\[
 \forall N\ \exists M\ge N:
 F(k,N,M)\text{ is unsatisfiable}.
\]
A finite cover with a threshold that changes with \(M\) therefore does not
imply a global assignment.
\end{remark}

\section*{Declarations}

\noindent\textbf{Generative AI contribution and proof provenance.}
\textbf{Substantive generative-AI use: this is a GPT-5.6-sol-generated
proof.} OpenAI's GPT-5.6-sol model, operating as Codex in a user-directed
research session, generated the central construction, the deterministic
semigroup sieve argument, the proof synthesis, and an initial manuscript
draft. Separate GPT-5.6-sol instances then performed line-by-line logical,
adversarial, and source-checking passes. These were internal AI-assisted
quality-control passes, not human peer review and not proof-assistant formal
verification. Defects identified in those passes were corrected in the
present manuscript.

GPT-5.6-sol is not listed as an author because a generative AI system cannot
assume scholarly accountability. By submitting this manuscript, the named
human author or authors certify that they have independently checked every
mathematical claim and bibliographic entry, approved the final manuscript,
and accept full responsibility for its correctness, originality, citations,
disclosure, and compliance with journal policy. Editors should apply the
AI-use policy of the target journal. This statement is intentionally more
detailed than a copy-editing disclosure because the AI system generated the
mathematical proof itself.

\medskip
\noindent\textbf{Human author responsibility.}
The corresponding author accepts responsibility for answering questions
about the argument and for investigating and resolving any concern about
accuracy or integrity.

\medskip
\noindent\textbf{Data, code, funding, and competing interests.}
No empirical data or finite computational result is used as a mathematical
premise. Every mathematical dependency is either proved in the text or
stated as a cited theorem with its hypotheses checked. No external funding
was used in the preparation of this manuscript. The author declares no
competing interests.

\medskip
\noindent\textbf{Acknowledgements.}
The proof-development process used separate AI model instances for
formal/local consistency checking, global/adversarial checking, and
verification of the cited analytic theorems. Their reports were internal
quality-control artifacts, not independent human referee reports, not
proof-assistant certificates, and not mathematical dependencies of the
proof.

\begin{thebibliography}{9}

\bibitem{BretècheTenenbaum}
R. de la Bret\`eche and G. Tenenbaum,
\emph{Remarks on the Selberg--Delange method},
Acta Arithmetica \textbf{200} (2021), no.~4, 349--369.

\bibitem{FouvryTenenbaum}
\'E. Fouvry and G. Tenenbaum,
\emph{Multiplicative functions in large arithmetic progressions and
applications},
Transactions of the American Mathematical Society \textbf{375} (2022),
no.~1, 245--299.

\bibitem{FriedlanderIwaniec}
J. Friedlander and H. Iwaniec,
\emph{Opera de Cribro},
American Mathematical Society Colloquium Publications, vol.~57,
American Mathematical Society, Providence, RI, 2010.

\bibitem{KurlbergEtAl}
P. Kurlberg, S. Lester, and L. Rosenzweig,
\emph{Superscars for arithmetic point scatters II},
Forum of Mathematics, Sigma \textbf{11} (2023), e42.

\end{thebibliography}

\end{document}
